\documentclass[activite]{thedocument}
\startdatakeys
\source{}
\BO{2026}
\cycle{4}
\competencesDuSocle{CA1,RA3}
\objectifsApprentissage{C4A1af,C4A1ah}
\motsCles{priorités opératoires, parenthèses}
\enddatakeys
\begin{document}
    \begin{enumerate}
        \item Adam a effectué les calculs suivants :
        $$3,5+7,5=11$$
        $$11\times8=88$$
        \'Elise regarde ses calculs et dit \guillemets{Je peux écrire tous tes
        calculs en une seule ligne~!} en lui montrant ceci~:
        $$3,5+7,5\times8$$
        A-t-elle raison~? Justifier.
        \item Calculer l'expression numérique suivante :$$4\times(8-5,5)$$
        \item \'Ecrire la règle permettant de calculer des expressions
        numériques contenants des parenthèses.
        \dotedline\dotedline\dotedline
    \end{enumerate}
\end{document}